% !TeX program = xelatex
% !TeX encoding = UTF-8
\documentclass[11pt,a4paper]{ctexart}
\usepackage[a4paper,margin=2.25cm]{geometry}
\usepackage{amsmath,amssymb,mathtools,bm}
\usepackage{booktabs,array}
\usepackage{graphicx}
\usepackage{xcolor}
\usepackage{tikz}
\usetikzlibrary{arrows.meta,positioning,fit,calc}
\usepackage{enumitem}
\usepackage{hyperref}
\hypersetup{colorlinks=true,linkcolor=black,urlcolor=blue,citecolor=black}
\setlength{\parindent}{2em}
\setlength{\parskip}{0.35em}
\setlist[itemize]{leftmargin=2em,itemsep=0.25em,topsep=0.25em}
\definecolor{mainblue}{RGB}{43,92,145}
\definecolor{lightblue}{RGB}{236,244,252}
\definecolor{lightgreen}{RGB}{236,248,240}
\definecolor{lightorange}{RGB}{253,244,232}
\definecolor{lightgray}{RGB}{246,246,246}

\title{\textbf{一个统一的组合优化算法评价框架}\\[0.35em]
\large 用 Optimal Transport 同时描述泛化性、性能与随机性}
\author{PitBench Methodology Note}
\date{2026 年 8 月}

\begin{document}
\maketitle

\begin{abstract}
本文给出一个简洁的统一视角：把组合优化算法看成一个带随机性的映射，它把实例分布映射成解或求解轨迹的分布。Optimal Transport（OT）的主要作用不是直接定义“泛化性”，而是为实例分布和输出分布提供可比较的几何结构。由此，\textbf{实例分布之间的 OT 距离}用于量化 distribution shift，\textbf{相对于最优解 anchor 的输出距离}用于量化性能，而\textbf{固定实例下输出分布的离散程度}用于量化算法随机性。三者因此可以在同一套概率与距离框架中描述。
\end{abstract}

\section{核心思想}

设组合优化实例属于空间 $\mathcal X$。算法记为 $A$。由于随机种子、并行调度、随机化启发式等因素，同一个实例上运行算法并不一定得到完全相同的求解过程，因此写成
\[
\tau = A(x;\xi),
\]
其中 $x\in\mathcal X$ 是实例，$\xi$ 表示算法内部随机性，$\tau$ 表示一次完整的求解轨迹（trajectory）。

我们再用一个特征映射
\[
\phi:\tau\mapsto y\in\mathcal Y
\]
把求解轨迹变成真正用于评价的对象。例如 $y$ 可以包含：
\[
\bigl(g(t_1),\ldots,g(t_K),\ \text{time-to-target},\ \text{nodes},\ \text{proof statistics}\bigr),
\]
其中 $g(t)$ 可以是相对于最优值的 gap 曲线。

因此，给定实例 $x$ 后，算法并不是产生一个确定的点，而是产生一个条件分布
\[
K_A(\mathrm dy\mid x).
\]
这个对象可以理解为：\textbf{固定实例 $x$，重复运行算法时，输出 $y$ 会如何分布。}

若实例本身来自某个经验总体 $P$，则算法在整个实例总体上诱导出输出分布
\[
Q_A^P(\mathrm dy)
=
\int_{\mathcal X} K_A(\mathrm dy\mid x)\,P(\mathrm dx).
\]
于是算法可以抽象为
\[
\boxed{P\xrightarrow{\quad A,\,\xi\quad}Q_A^P.}
\]
这就是统一框架的基本对象。

\begin{figure}[htbp]
\centering
\begin{tikzpicture}[>=Latex,node distance=2.0cm]
\node[draw,rounded corners,fill=lightblue,minimum width=3.2cm,minimum height=1.05cm] (P) {实例分布 $P$};
\node[draw,rounded corners,fill=lightgreen,right=3.6cm of P,minimum width=3.6cm,minimum height=1.05cm] (Q) {输出分布 $Q_A^P$};
\draw[->,very thick,mainblue] (P) -- node[above]{随机算法 $A(\cdot;\xi)$} (Q);
\node[below=0.45cm of P,align=center] {经验实例总体\\$\sum_i a_i\delta_{x_i}$};
\node[below=0.45cm of Q,align=center] {解 / 轨迹 / 性能特征\\的概率分布};
\end{tikzpicture}
\caption{算法被视为从实例总体到输出总体的随机映射。}
\end{figure}

\section{为什么需要 OT：给“变化”一个距离}

真实 benchmark 中，我们通常只有有限实例
\[
x_1,\ldots,x_n,
\]
并不需要假设它们来自高斯分布，也不需要假设实例空间形成光滑流形。最自然的表示是经验测度
\[
\widehat P=\sum_{i=1}^n a_i\delta_{x_i},
\qquad a_i\ge 0,\quad \sum_i a_i=1.
\]
另一组实例同样写成
\[
\widehat Q=\sum_{j=1}^m b_j\delta_{x'_j}.
\]

如果已经定义了单个实例之间的 ground distance
\[
d_{\mathcal X}(x,x'),
\]
那么 $p$-Wasserstein 距离定义为
\[
W_p(\widehat P,\widehat Q)
=
\left[
\min_{\pi\in\Pi(a,b)}
\sum_{i=1}^n\sum_{j=1}^m
\pi_{ij}\,d_{\mathcal X}(x_i,x'_j)^p
\right]^{1/p},
\]
其中 $\pi_{ij}$ 表示把多少概率质量从 $x_i$ 搬到 $x'_j$，并满足
\[
\sum_j\pi_{ij}=a_i,
\qquad
\sum_i\pi_{ij}=b_j,
\qquad
\pi_{ij}\ge 0.
\]

因此 OT 的语义非常直接：
\[
\boxed{\text{把一个实例总体变成另一个实例总体，最少需要搬运多少“概率质量 $\times$ 距离”？}}
\]

OT 在这里提供的是一个\textbf{distribution-shift geometry}。它本身不是“泛化性”，而是告诉我们实例分布到底改变了多少。

\section{三件事如何进入同一个框架}

\subsection{泛化性：实例分布改变后，算法性能恶化多少？}

设 $P$ 是开发/训练时的实例总体，$Q$ 是另一个测试总体。输入侧的变化量由
\[
W_{\mathcal X}(P,Q)
\]
描述。

为了评价算法本身，需要先定义一个相对于最优解的 loss。设实例 $x$ 的最优解或最优行为对应 anchor
\[
y^*(x).
\]
给定距离 $d_{\mathcal Y}$，定义算法在总体 $P$ 上的 oracle-relative risk：
\[
L_A(P)
=
\mathbb E_{x\sim P}
\mathbb E_{Y\sim K_A(\cdot\mid x)}
\bigl[d_{\mathcal Y}(Y,y^*(x))\bigr].
\]

那么一个最简单的 distributional sensitivity 是
\[
\boxed{
R_A(P,Q)
=
\frac{|L_A(P)-L_A(Q)|}{W_{\mathcal X}(P,Q)}
}
\]
（当分母非零时）。它表达：\textbf{实例总体移动单位距离时，算法相对于最优解的性能变化有多大。}

这里有一个重要原则：
\[
W_{\mathcal Y}(Q_A^P,Q_A^Q)\ \text{小}
\]
\textbf{并不自动意味着泛化好。} 因为实例变化以后，真正的最优解结构本身也可能应该发生很大变化。泛化性应关注的是\textbf{oracle-relative performance 是否稳定}，而不是要求输出“不变化”。

\subsection{性能：算法离最优 anchor 有多远？}

固定一个实例 $x$。算法输出是分布 $K_A(\cdot\mid x)$，最优解 anchor 可以写成一个 Dirac 测度
\[
\delta_{y^*(x)}.
\]
于是可以定义
\[
\boxed{
\mathrm{Perf}_A(x)
=
W_{\mathcal Y}
\bigl(K_A(\cdot\mid x),\delta_{y^*(x)}\bigr)
}
\]
或更直接地使用期望距离
\[
\mathrm{Perf}_A(x)
=
\mathbb E_{Y\sim K_A(\cdot\mid x)}
[d_{\mathcal Y}(Y,y^*(x))].
\]
越小表示算法越接近最优。

对于 anytime solver，可以把输出定义为整条最优性 gap 曲线，例如
\[
g_x(t)
=
\frac{f_x(t)-f^*(x)}{s(x)},
\]
其中 $f^*(x)$ 是最优目标值，$s(x)$ 是归一化尺度。此时 $g_x(t)=0$ 就是统一的最优 anchor。

\subsection{随机性：固定实例时，输出分布有多分散？}

随机性必须在\textbf{固定实例}后测量，否则实例异质性会和算法随机性混在一起。

给定 $x$，独立运行两次算法：
\[
Y,Y'\overset{\mathrm{iid}}{\sim}K_A(\cdot\mid x).
\]
一个简单的随机性/不稳定性指标是
\[
\boxed{
S_A(x)
=
\frac12
\mathbb E\bigl[d_{\mathcal Y}(Y,Y')^2\bigr].
}
\]
若 $S_A(x)$ 很小，则不同 seed 得到的行为相近；若很大，则算法对 seed 或其他随机源高度敏感。

在一个实例总体 $P$ 上，可以进一步聚合：
\[
S_A(P)=\mathbb E_{x\sim P}[S_A(x)].
\]

因此，随机性并不是额外附加的指标，而是由同一个条件输出分布
\[
K_A(\cdot\mid x)
\]
自然给出。

\section{统一图景}

三类评价可以放进同一个图中：

\begin{figure}[htbp]
\centering
\resizebox{0.96\textwidth}{!}{%
\begin{tikzpicture}[>=Latex,node distance=2.0cm]
\node[draw,rounded corners,fill=lightblue,minimum width=3.2cm,minimum height=1cm] (P) {实例总体 $P$};
\node[draw,rounded corners,fill=lightblue,below=2.2cm of P,minimum width=3.2cm,minimum height=1cm] (Q) {实例总体 $Q$};
\node[draw,rounded corners,fill=lightgreen,right=4.3cm of P,minimum width=3.5cm,minimum height=1cm] (AP) {$Q_A^P$};
\node[draw,rounded corners,fill=lightgreen,right=4.3cm of Q,minimum width=3.5cm,minimum height=1cm] (AQ) {$Q_A^Q$};
\node[draw,rounded corners,fill=lightorange,right=3.8cm of AP,minimum width=3.1cm,minimum height=1cm] (OP) {oracle $Q_*^P$};
\node[draw,rounded corners,fill=lightorange,right=3.8cm of AQ,minimum width=3.1cm,minimum height=1cm] (OQ) {oracle $Q_*^Q$};

\draw[->,very thick,mainblue] (P) -- node[above]{$A,\xi$} (AP);
\draw[->,very thick,mainblue] (Q) -- node[above]{$A,\xi$} (AQ);
\draw[<->,thick] (P) -- node[left,align=center]{$W_{\mathcal X}(P,Q)$\\实例分布变化} (Q);
\draw[<->,thick] (AP) -- node[above]{performance} (OP);
\draw[<->,thick] (AQ) -- node[above]{performance} (OQ);
\draw[<->,dashed,thick] (AP) -- node[right,align=center]{输出变化\\（描述性）} (AQ);
\end{tikzpicture}%
}
\caption{统一框架：输入侧 OT 描述 distribution shift；输出相对 oracle 的距离描述性能；固定实例下输出条件分布的 dispersion 描述随机性。}
\end{figure}

因此三件事可以概括为：
\[
\boxed{
\begin{aligned}
\text{Performance: }&\quad
\mathbb E_{x\sim P}W_{\mathcal Y}
\bigl(K_A(\cdot\mid x),\delta_{y^*(x)}\bigr),\\[0.4em]
\text{Randomness: }&\quad
\mathbb E_{x\sim P}\frac12\mathbb E[d_{\mathcal Y}(Y,Y')^2\mid x],\\[0.4em]
\text{Generalization: }&\quad
\text{oracle-relative performance 随 }W_{\mathcal X}(P,Q)\text{ 的变化规律。}
\end{aligned}}
\]

\section{OT 在输入侧与输出侧分别做什么}

为了避免概念混淆，可以区分两种 OT：

\begin{center}
\begin{tabular}{>{\raggedright\arraybackslash}p{3.2cm} >{\raggedright\arraybackslash}p{5.1cm} >{\raggedright\arraybackslash}p{5.1cm}}
\toprule
 & 输入侧 OT & 输出侧 OT \\
\midrule
比较对象 & 两个实例或实例总体 & 两个解 / 轨迹 / 输出分布 \\
核心问题 & instance distribution 改变了多少？ & solver behavior 改变了多少？ \\
主要用途 & 泛化、distribution shift、DRO & 性能差异、行为差异、随机性描述 \\
关键前提 & instance ground metric 有意义 & output/trajectory representation 有意义 \\
\bottomrule
\end{tabular}
\end{center}

其中，单个复杂实例内部是否需要使用所谓 Inner OT，是一个\textbf{可选的 ground-metric 构造方法}，不是统一框架成立的前提。Outer OT 才直接对应“两个实例总体之间的距离”。

\section{与 PitBench 的关系：评价 patch 的泛化}

如果 $B$ 是 base solver，$A$ 是经过 agent 或 human patch 后的 solver，可以定义 patch 在实例 $x$ 上的期望收益
\[
\Delta_{A/B}(x)
=
L_B(x)-L_A(x).
\]
总体 $P$ 上的平均收益为
\[
G_{A/B}(P)
=
\mathbb E_{x\sim P}[\Delta_{A/B}(x)].
\]
于是 patch 的泛化问题变成：
\[
\boxed{
W_{\mathcal X}(P,Q)\ \text{增大时，}\ G_{A/B}(Q)\ \text{还能保留多少？}
}
\]
这比只报告一个 dev/test speedup 更丰富，因为它试图回答：\textbf{patch 的收益随着实例分布偏移究竟如何衰减。}

如果未来 ground geometry 被可靠验证，还可以进一步定义 Wasserstein ball
\[
\mathcal U_\varepsilon(P)
=
\{Q:W_{\mathcal X}(Q,P)\le \varepsilon\}
\]
以及最坏情况下的 retained gain
\[
G_{A/B}^{\mathrm{rob}}(\varepsilon)
=
\inf_{Q\in\mathcal U_\varepsilon(P)}G_{A/B}(Q).
\]
这时 $\varepsilon$ 就有明确含义：允许实例总体在所定义的 geometry 中移动多远。

\section{最重要的边界条件}

这套框架的价值在于统一，而不是强迫所有对象都使用 OT。需要保留以下边界：

\begin{itemize}
\item OT 不自动给出正确的 ground metric；实例表示如果与 solver response 无关，Wasserstein 距离也没有评价意义。
\item 泛化性不是“输出分布变化越小越好”。实例改变时，正确的最优解本身也可能必须显著改变。
\item 随机性必须条件于固定实例测量，否则会把 instance heterogeneity 错当成 solver stochasticity。
\item optimal anchor 描述的是 problem-level quality；human patch 则是 software-engineering reference，两者含义不同。
\item 对真实 benchmark，应优先使用经验测度，不需要先假设分布可参数化，也不需要假设实例空间形成光滑流形。
\end{itemize}

\section{一句话总结}

\begin{center}
\fcolorbox{mainblue}{lightblue}{%
\begin{minipage}{0.91\textwidth}
\centering
\textbf{把算法视为从经验实例分布到随机解/轨迹分布的 stochastic mapping。}\\[0.4em]
输入侧 OT 给出 distribution shift 的 geometry；最优解 anchor 给出绝对性能参照；固定实例下的条件输出分布给出算法随机性。由此，performance、stability 和 distributional generalization 可以在同一个概率与距离框架中被统一描述。
\end{minipage}}
\end{center}

\end{document}
